\section{Herramientas y automatización}

\subsection{Pipeline de automatización}

Offline RL necesita encadenar data -> training -> eval -> deploy en un pipeline. Componentes típicos:

\begin{itemize}
    \item Data versioning: cada entrenamiento apunta a los datos, el reward log y el ID de la behavior policy.
    \item Training orchestration: Pipeline.com, Kubeflow pipeline disparan el training job.
    \item Evaluation orchestration: toma de manera automatizada key states del replay buffer y corre el benchmark.
    \item Deployment gate: si las metrics son anómalas, revierte automáticamente y notifica al owner.
\end{itemize}

\begin{warningbox}{Una interrupción del pipeline provoca drift}
En cuanto falta cualquier eslabón del pipeline (data/eval/deploy), los cambios posteriores dejan de poder rastrearse. Es necesario exigir que todo code change pase por un artifact generado por el pipeline, para evitar «correrlo manualmente».
\end{warningbox}

\subsection{Herramientas de apoyo e integración}

El docente enumeró varias herramientas que se integran bien con offline RL:

\begin{itemize}
    \item \textbf{Data-versioned storage} (Feathr, Deep Lake) almacena el metadata del dataset.
    \item \textbf{Experiment tracking} (Weights \& Biases) registra la Q variance y el KL drift.
    \item \textbf{Policy registry}: registra el policy hash, el deploy timestamp, el sign-off y el enlace al governance doc.
    \item \textbf{Replay simulator}: hace un replay de las trayectorias de entrenamiento con deterministic seeds y registra las divergences.
\end{itemize}

\subsection{Resumen de la sección}

El automation pipeline no es solo un pipeline: convierte cada artifact de offline RL en un producto rastreable. El logging en tiempo real más la integración de herramientas facilitan la gobernanza.
